\clearpage
\fancyhead[L]{\small BEOBACHTUNG VON ABLENKUNGEN.}
\SourceLeaf{401}
\pdfbookmark[0]{Über ein Mittel, die Beobachtung von Ablenkungen zu erleichtern. 1839}{ablenkungen-1839}
\begin{center}
{\large ÜBER EIN MITTEL,\par\medskip DIE BEOBACHTUNG VON ABLENKUNGEN ZU ERLEICHTERN.}
\par\medskip
Resultate aus den Beobachtungen des magnetischen Vereins. 1839. II.
\end{center}
\par\medskip\Needspace{4\baselineskip}\pdfbookmark[1]{1.}{ablenkungen-art-1}\centerline{\textbf{1.}}\nobreak\smallskip
Tritt zur erdmagnetischen Kraft eine weitere Kraft hinzu, die stetig auf die Nadel eines Magnetometers wirkt, aber in einer gegen den magnetischen Meridian geneigten Richtung, so nimmt die Nadel eine neue Gleichgewichtsstellung ein. Die Größe der Ablenkung kann zur Messung der Zusatzkraft dienen. Damit sich die Ablenkung messen lässt, muss die neue Gleichgewichtsstellung noch innerhalb der Skala liegen. Kann oder will man den völligen Ruhezustand der Nadel nicht abwarten, dürfen außerdem die noch stattfindenden Schwingungen die Grenzen der Skala nicht überschreiten. War die Nadel in Ruhe, solange der Erdmagnetismus allein auf sie wirkte, und lässt man die Zusatzkraft auf einmal voll wirken, so beginnt die Nadel eine Schwingung, deren Mitte die neue Gleichgewichtsstellung ist. Die vorige Stellung bildet den einen Elongationspunkt, der zweite liegt ebenso weit von der Mitte entfernt auf der gegenüberliegenden Seite. Liegt die neue Gleichgewichtsstellung zwar innerhalb der Skala, aber nahe an ihrer Grenze, so müsste man bei der langsamen Abnahme des Schwingungsbogens ohne künstliche Hilfsmittel lange warten, bis sich dieser Punkt auf diese Weise bestimmen ließe. Dadurch würden aber in allen Fällen schon wegen der stündlichen Änderung der Deklination die Zuverlässigkeit und Brauchbarkeit der Bestimmung stark beeinträchtigt. Fast ganz vereitelt würde die Bestimmung
\expandafter\def\csname HistoricalNoteText404\endcsname{Resultate, 1837. IV. [S. 383. d. B.]}
\renewcommand{\Article}[1]{\par\medskip\Needspace{4\baselineskip}\pdfbookmark[1]{#1.}{ablenkungen-art-#1}\centerline{\textbf{#1.}}\nobreak\smallskip}%
\SourceLeaf{402}in Fällen, in denen sich die Stärke der Zusatzkraft schon in kurzer Zeit beträchtlich ändert, wie bei galvanischen Strömen.

\Article{2}
Das folgende einfache Verfahren behebt diesen Nachteil. Man lässt die Zusatzkraft zunächst nur während eines Drittels der Schwingungsdauer wirken. Dann unterbricht man ihre Wirkung für eine ebenso lange Zeit und lässt sie erst danach dauerhaft wirken. Beträgt also etwa die Schwingungsdauer der Magnetometernadel 30 Sekunden und soll die durch einen galvanischen Strom erzeugte Ablenkung gemessen werden, so schließt man den Stromkreis bei einem Sekundenschlag, den man als 0 zählt, öffnet ihn bei $10''$ wieder und schließt ihn schließlich bei $20''$ endgültig. Soll die Ablenkung durch einen Magnetstab erzeugt werden, den man an eine bestimmte Stelle legt, so nähert man sich der zuvor genau und zweckmäßig markierten Stelle mit zunächst senkrecht gehaltenem Magnetstab. Bei $0''$ legt man ihn plötzlich nieder, bei $10''$ richtet man ihn ebenso schnell wieder auf, und bei $20''$ legt man ihn zum zweiten Mal endgültig hin. Das Ergebnis ist, dass sich die Nadel während der ersten 10 Sekunden aus ihrer ursprünglichen Ruhestellung mit zunehmender Geschwindigkeit der Stellung nähert, die der Ablenkung entspricht. Bei $10''$ hat sie genau die Mitte zwischen beiden Stellungen erreicht. Während der zweiten 10 Sekunden durchläuft sie dann die andere Hälfte der Strecke mit verzögerter Bewegung, sodass sie bei 20 Sekunden die neue Stellung erreicht und jede Bewegung verloren haben wird.

Man sieht leicht, dass man die Nadel auf ganz ähnliche Weise von einer ruhenden abgelenkten Stellung in die entgegengesetzte überführen kann, sodass sie dort ohne Bewegung ist. Man lässt dazu die ablenkende Kraft während eines Drittels der Schwingungsdauer in der entgegengesetzten Richtung wirken, dann ebenso lange wieder in der früheren Richtung und wechselt darauf erneut. Bei galvanischen Strömen lässt sich die Richtung mit einem geeigneten Kommutator fast augenblicklich wechseln. Bei ablenkenden Magnetstäben geschieht das durch eine rasche halbe Umdrehung, am bequemsten in der Horizontalen, sodass der Nordpol des Stabs an die Stelle des Südpols kommt.

Schließlich ist klar, dass man die Nadel nach Beobachtung der Ablenkung auf dieselbe Weise wieder im reinen magnetischen Meridian zur Ruhe bringen kann. Man braucht die ablenkende Kraft nur zuerst während eines Drittels der Schwin\SourceLeaf{403}gungsdauer auszusetzen, sie dann noch einmal ebenso lange wirken und schließlich aufhören zu lassen.

\Article{3}
Das beschriebene Verfahren beruht auf der Annahme dass

\textit{erstens} die Nadel so schwingt, dass ihr Abstand von der Mitte der Schwingung, solange diese Mitte selbst nicht verändert wird, dem Sinus eines Winkels proportional bleibt, der sich gleichförmig ändert und während einer Schwingungsdauer um $180^0$ zunimmt, und

\textit{zweitens}, die Zusatzkraft die Schwingungsdauer nicht verändert.

Soweit beide Annahmen nicht \textit{mit absoluter Genauigkeit} gelten und bei der Ausführung weder die Wechsel ganz augenblicklich geschehen noch die vorgeschriebenen Zwischenzeiten absolut genau eingehalten werden können, wird man die Nadel am Ende des Vorgangs tatsächlich selten in vollkommener Ruhe antreffen. Für den praktischen Zweck genügt es aber schon, wenn die verbleibende Bewegung so gering ist, dass man sofort auf die übliche Weise mit der Beobachtung der wahren Gleichgewichtsstellung beginnen kann.

Unter den gegebenen Umständen können diese Annahmen nur sehr wenig von der Wahrheit abweichen. Schon die Anwendbarkeit des Magnetometers setzt voraus, dass die Zusatzkraft nur eine mäßige Ablenkung hervorruft. Dabei gilt das in der ersten Annahme enthaltene Gesetz hinreichend genau, abgesehen von einem Ausnahmefall, der gleich zur Sprache kommt. Die Änderung der Schwingungsdauer durch die ablenkende Kraft ist ganz unmerklich, wenn diese senkrecht zum magnetischen Meridian wirkt, wie es fast immer der Fall ist. Wirkte sie aber auch in einer schrägen Richtung, so bliebe die dadurch bewirkte Änderung der Schwingungsdauer, sofern die Kraft selbst nur ein kleiner Bruchteil der erdmagnetischen Kraft ist, für die kurze Zeit des Vorgangs doch ganz unerheblich.

Nur ein Fall ist auszunehmen, nämlich wenn die Nadel unter dem Einfluss eines Dämpfers schwingt, der den Schwingungsbogen beträchtlich verkleinert. Dann gilt das obige Gesetz nicht mehr, und die genaue Befolgung des beschriebenen Verfahrens würde nicht zum Ziel führen. Auf der anderen Seite ist dann allerdings auch der
\SourceLeaf{404}in Artikel 1 erwähnte Nachteil viel geringer, da ein kräftiger Dämpfer die Nadel von selbst in mäßiger Zeit zur Ruhe bringt. Da das Verfahren in diesem Fall aber nur einer Abänderung bedarf, um denselben Erfolg zu erzielen, und man jeden unnötigen Zeitverlust stets gern vermeiden möchte, lohnt es sich praktisch wie theoretisch. die Frage ganz allgemein zu betrachten.

\Article{4}
Zunächst haben wir folgende allgemeine Aufgabe zu lösen.

Ein Magnetstab schwingt, während die auf ihn wirkenden Kräfte wiederholt verändert werden. Dabei bleiben jedoch die Schwingungsdauer und das logarithmische Dekrement\HistoricalNote{404} unverändert und die Schwingungsbögen klein genug, um Größen dritter Ordnung vernachlässigen zu können. Aus dem anfänglichen Bewegungszustand soll derjenige abgeleitet werden, der nach der letzten Änderung eintritt.

Es sei $T$ die Schwingungsdauer, $\epsilon$ das logarithmische Dekrement, $e$ die Basis der hyperbolischen Logarithmen, $m$ der Modul der Briggsschen Logarithmen und $\pi$ das Verhältnis des Kreisumfangs zum Durchmesser. Man setze
\[
n=\frac{\pi}{T},\qquad \epsilon=\frac{\lambda}{mT}
\]
Unter den obigen Voraussetzungen wird demnach die Stellung $x$ zur Zeit $t$ durch die Formel ausgedrückt
\[
x=p+Ae^{-\epsilon t}\sin(nt-B)
\]
der man auch die Form geben kann
\[
x=p+ae^{-\epsilon t}\cos nt+be^{-\epsilon t}\sin nt
\]
wobei $p$ die Gleichgewichtsstellung bezeichnet und die Koeffizienten $a,b$ so lange konstant bleiben, wie $p$ konstant ist. Daraus ergibt sich die Geschwindigkeit zu
\[
\frac{dx}{dt}=-e^{-\epsilon t}(na\sin nt+\epsilon a\cos nt-nb\cos nt+\epsilon b\sin nt)
\]
oder, wenn man einen Hilfswinkel $\varphi$ so einführt, dass $\frac{\epsilon}{n}=\operatorname{tang}\varphi$ gilt,
\SourceLeaf{405}\[
\frac{dx}{dt}=-\frac{ne^{-\epsilon t}}{\cos\varphi}(a\sin(nt+\varphi)-b\cos(nt+\varphi))
\]
Für $ae^{-\epsilon t}\cos nt+be^{-\epsilon t}\sin nt$ schreiben wir $u$, sodass $x=p+u$ gilt.

Nun seien $t',t'',t'''$ die bestimmten Werte von $t$, zu denen die wirkende Kraft geändert wird. Ferner seien die bestimmten Werte von $p,a,b$ in den verschiedenen Zeitabschnitten die folgenden
\[
\begin{array}{lll}
p^0,\ a^0,\ b^0&\text{vor}&t'\\
p',\ a',\ b'&\text{von}&t'\ \text{bis}\ t''\\
p'',\ a'',\ b''&\text{von}&t''\ \text{bis}\ t'''\\
p''',\ a''',\ b'''&\text{nach}&t'''
\end{array}
\]
Schließlich gehe der allgemeine Ausdruck für $u$, wenn man für $a$ und $b$ die bestimmten Werte einsetzt, in $u^0,u',u'',u'''$ über, sodass vor dem ersten Wechsel $x=p^0+u^0$ gilt, von da bis zum zweiten $x=p'+u'$ und so weiter.

Da der Augenblick $t'$ zugleich der letzte des ersten Zeitabschnitts und der erste des folgenden ist, müssen für $t=t'$ sowohl $x$ als auch $\frac{dx}{dt}$ denselben Wert erhalten, ob man nun in den obigen allgemeinen Ausdrücken für $p,a,b$ die Werte $p^0,a^0,b^0$ oder $p',a',b'$ einsetzt.

Also gilt
\[
\begin{aligned}
0&=p'-p^0+(a'-a^0)e^{-\epsilon t'}\cos nt'+(b'-b^0)e^{-\epsilon t'}\sin nt'\\
0&=(a'-a^0)\sin(nt'+\varphi)-(b'-b^0)\cos(nt'+\varphi)
\end{aligned}
\]
woraus man leicht ableitet
\[
\begin{aligned}
a'-a^0&=-\frac{p'-p^0}{\cos\varphi}.e^{\epsilon t'}\cos(nt'+\varphi)\\
b'-b^0&=-\frac{p'-p^0}{\cos\varphi}.e^{\epsilon t'}\sin(nt'+\varphi)
\end{aligned}
\]
und daraus
\[
u'=u^0-\frac{p'-p^0}{\cos\varphi}.e^{-\epsilon(t-t')}\cos(n(t-t')-\varphi)
\]
Auf dieselbe Weise erhält man
\[
\begin{aligned}
u''&=u'-\frac{p''-p'}{\cos\varphi}.e^{-\epsilon(t-t'')}\cos(n(t-t'')-\varphi)\\
u'''&=u''-\frac{p'''-p''}{\cos\varphi}.e^{-\epsilon(t-t''')}\cos(n(t-t''')-\varphi)
\end{aligned}
\]
und so weiter, wenn sich die bewegenden Kräfte noch öfter ändern.

\SourceLeaf{406}So bestimmt der anfängliche Bewegungszustand jeden nachfolgenden.

\Article{5}
Für den hier untersuchten Fall ist $p''=p^0$ und $p'''=p'$ zu setzen. Dadurch erhält man
\[
\begin{aligned}
u'''={}&u^0-\frac{p'-p^0}{\cos\varphi}.e^{-\epsilon t}.\\
&.\{e^{\epsilon t'}\cos(n(t-t')-\varphi)-e^{\epsilon t''}\cos(n(t-t'')-\varphi)\\
&\qquad+e^{\epsilon t'''}\cos(n(t-t''')-\varphi)\}
\end{aligned}
\]
Setzt man
\[
\begin{aligned}
e^{-\epsilon(t''-t')}\cos n(t''-t')-1+e^{\epsilon(t'''-t'')}\cos n(t'''-t'')&=f\\
e^{-\epsilon(t''-t')}\sin n(t''-t')\phantom{-1}{}-e^{\epsilon(t'''-t'')}\sin n(t'''-t'')&=g
\end{aligned}
\]
so geht diese Formel über in
\[
\begin{aligned}
u'''={}&u^0-\frac{p'-p^0}{\cos\varphi}.e^{-\epsilon(t-t'')}.\{f\cos(n(t-t'')-\varphi)\\
&\qquad-g\sin(n(t-t'')-\varphi)\}
\end{aligned}
\]
Daraus folgt, dass, wenn die Zwischenzeiten $t''-t',t'''-t''$ so bestimmt sind, dass $f=0$ und $g=0$ gilt, allgemein
\[
u'''=u^0
\]
oder $a'''=a^0,b'''=b^0$ gilt.

War also die Nadel vor den Wechseln in $p^0$ in Ruhe, so wird sie nach ihnen in $p'$ in Ruhe sein. Im entgegengesetzten Fall hat die Nadel nach den drei Wechseln in jedem Augenblick genau dieselbe Geschwindigkeit und dieselbe Stellung zum Mittelpunkt ihrer Bewegung $p'$, die sie im selben Augenblick relativ zu $p^0$ hätte, wenn sie ihre ursprüngliche Bewegung ungestört fortgesetzt hätte. Mit einem Wort, nur der Mittelpunkt der Bewegung wird versetzt, und die Bewegung selbst wird überhaupt nicht verändert.

\Article{6}
Es bleibt noch, die Zwischenzeiten so zu bestimmen, dass die Gleichungen $f=0,g=0$ erfüllt sind. Setzt man
\SourceLeaf{407}\[
t''-t'=qT,\qquad t'''-t''=rT
\]
und erinnert sich daran, dass $e=10^m$, so lauten diese Gleichungen
\[
\begin{aligned}
10^{-q\lambda}\cos q\pi+10^{r\lambda}\cos r\pi&=1\\
10^{-q\lambda}\sin q\pi&=10^{r\lambda}\sin r\pi
\end{aligned}
\]
Nimmt der Schwingungsbogen unmerklich ab, so muss man also $\cos q\pi+\cos r\pi=1$ und $\sin q\pi=\sin r\pi$ setzen, und damit
\[
q\pi=r\pi=60^0\ \text{oder}\ =\tfrac{1}{3}\pi,\quad\text{und}\quad t''-t'=t'''-t''=\tfrac{1}{3}T
\]
wie schon in Artikel 2 bemerkt wurde. Bei einem merklichen logarithmischen Dekrement sind diese Gleichungen dagegen auf indirektem Weg zu lösen. Man kann der Rechnung die folgende Form geben.

Aus der Verbindung der Gleichungen folgt
\[
\begin{gathered}
\operatorname{tang}r\pi=\frac{\sin q\pi}{10^{q\lambda}-\cos q\pi},\qquad
10^{2r\lambda}=1-2.10^{-q\lambda}\cos q\pi+10^{-2q\lambda}
\end{gathered}
\]
Durch Elimination von $r$ erhält man also die Gleichung mit einer unbekannten Größe
\[
\frac{\pi}{2\lambda}\log(1-2.10^{-q\lambda}\cos q\pi+10^{-2q\lambda})=\operatorname{Arc.\,tang}\frac{\sin q\pi}{10^{q\lambda}-\cos q\pi}
\]
wobei der Briggssche Logarithmus gemeint ist. Ist diese Gleichung erfüllt, so hat man offenbar zugleich den Wert von $r$.

\Article{7}
Um denen, die das beschriebene Verfahren mit einem Dämpfer ausführen wollen, die im vorigen Artikel erklärte Rechnung zu ersparen, teile ich hier eine im Voraus berechnete Tafel mit. Aus ihr lässt sich für jedes logarithmische Dekrement das Verhältnis der beiden Zwischenzeiten zur Schwingungsdauer sofort ablesen. Man sieht daraus, dass mit zunehmendem logarithmischem Dekrement die erste Zwischenzeit immer größer, die zweite immer kleiner wird. Ihre Summe ist zwar nur für $\lambda=0$ genau gleich zwei Dritteln der Schwingungsdauer, weicht aber viel langsamer davon ab. Dass es bei der praktischen Anwendung genügt, etwa nur die ersten Dezimalstellen der Werte von $q$ und $r$ zu berücksichtigen, braucht kaum erwähnt zu werden.

\Needspace{400pt}\SourceLeaf{408}\begin{center}\textit{\textbf{Tafel.}}\end{center}
\begingroup\fontsize{9}{11}\selectfont
\begin{center}
\begin{tabular}{c|c|c||c|c|c}
$\lambda$&$q$&$r$&$\lambda$&$q$&$r$\\\hline
0&0.33333&0.33333&0.30&0.45921&0.21406\\
0.01&0.33757&0.32911&0.31&0.46322&0.21048\\
0.02&0.34181&0.32489&0.32&0.46721&0.20694\\
0.03&0.34606&0.32068&0.33&0.47118&0.20343\\
0.04&0.35031&0.31648&0.34&0.47513&0.19996\\
0.05&0.35456&0.31229&0.35&0.47906&0.19652\\[3pt]
0.06&0.35882&0.30812&0.36&0.48297&0.19311\\
0.07&0.36308&0.30395&0.37&0.48685&0.18975\\
0.08&0.36734&0.29981&0.38&0.49071&0.18641\\
0.09&0.37160&0.29568&0.39&0.49454&0.18311\\
0.10&0.37585&0.29156&0.40&0.49835&0.17985\\[3pt]
0.11&0.38011&0.28746&0.41&0.50214&0.17663\\
0.12&0.38436&0.28338&0.42&0.50590&0.17344\\
0.13&0.38861&0.27932&0.43&0.50963&0.17029\\
0.14&0.39285&0.27528&0.44&0.51334&0.16718\\
0.15&0.39708&0.27126&0.45&0.51702&0.16411\\[3pt]
0.16&0.40131&0.26727&0.46&0.52067&0.16107\\
0.17&0.40552&0.26329&0.47&0.52430&0.15808\\
0.18&0.40973&0.25934&0.48&0.52790&0.15512\\
0.19&0.41393&0.25542&0.49&0.53147&0.15220\\
0.20&0.41812&0.25152&0.50&0.53501&0.14931\\[3pt]
0.21&0.42230&0.24764&0.51&0.53852&0.14647\\
0.22&0.42646&0.24379&0.52&0.54201&0.14367\\
0.23&0.43061&0.23997&0.53&0.54546&0.14090\\
0.24&0.43474&0.23618&0.54&0.54889&0.13817\\
0.25&0.43886&0.23242&0.55&0.55229&0.13548\\[3pt]
0.26&0.44297&0.22868&0.56&0.55566&0.13283\\
0.27&0.44705&0.22498&0.57&0.55900&0.13022\\
0.28&0.45112&0.22131&0.58&0.56231&0.12765\\
0.29&0.45517&0.21767&0.59&0.56559&0.12511\\
0.30&0.45921&0.21406&0.60&0.56884&0.12261
\end{tabular}
\end{center}\endgroup

\Article{8}
Die unserer Theorie zugrunde liegende Annahme, dass die drei Wechsel augenblicklich geschehen, trifft bei der tatsächlichen Ausführung des Verfahrens niemals in aller Genauigkeit zu, obwohl man bei Ablenkungen durch galvanische Ströme die für jeden Wechsel nötige Zeit als unmerklich betrachten kann. Bei Ablenkungen durch Magnetstäbe ist diese Zeit dagegen je nach Größe und Gewicht der Stäbe mehr oder weniger beträchtlich. Bei fünfundzwanzigpfündigen Stäben braucht man für einen Wechsel stets mehrere Sekunden, besonders wenn es um einen Wechsel zwischen zwei entgegengesetzten La\SourceLeaf{409}gen geht und nicht um einen zwischen senkrechter und waagerechter Lage. Für diesen Fall, der tatsächlich bei Weitem der wichtigste und häufigste ist, lässt sich der Vorgang aber leicht so einrichten, dass das Ergebnis kaum merklich beeinträchtigt wird. Man muss nur darauf achten, dass der zweite und dritte Wechsel auf dieselbe Weise wie der erste geschehen, also auch gleich lange dauern, und dass man diese Zeit von den sonst nötigen Zwischenzeiten abzieht. Beträgt zum Beispiel (wie in \textit{Res.} 1837. IV.) [S. 390. d. B.] das logarithmische Dekrement 0.33570 und die Schwingungsdauer $21''21439$, so ergibt die obige Tafel für die erste Zwischenzeit $=10''04$ und für die zweite $=4''27$. Stellt sich nun heraus, dass man zur Ausführung eines Wechsels drei Sekunden braucht, so \textit{beginnt} man den ersten Wechsel bei $0''$. Von $3''$ bis $10''$. bleibt der Stab in der neuen Lage. Durch den bei $10''$ \textit{beginnenden} neuen Wechsel ist der Stab bei $13''$ in die erste Lage zurückgebracht. Dort bleibt er nur $1\tfrac{1}{4}$ Sekunden liegen, worauf der dritte Wechsel beginnt, sodass der ganze Vorgang erst bei $17\tfrac{1}{4}$ Sekunden abgeschlossen ist. Eine ausführlichere Untersuchung, die hier aus Platzgründen übergangen werden muss, ergibt nämlich, dass, wenn $p^0$ in $p'$ allmählich übergeht statt sprunghaft, ebenso beim zweiten Wechsel $p'$ in $p^0$ und beim dritten wieder $p^0$ in $p'$, das Ergebnis ganz so bleibt, wie es am Schluss von Artikel 5 angegeben ist. Dazu müssen nur die drei Übergangszeiten gleich lang sein, die drei Übergänge selbst in ähnlichen Stufenfolgen verlaufen und die berechneten Zwischenzeiten $qT,rT$ auf die Anfangszeitpunkte der Wechsel bezogen oder, was dasselbe ist, die beiden ersten Übergangszeiten in sie eingerechnet werden.

